% =============================================================
%  Chapter 2 — Background Review : BRIEF SKELETON
%  Just the structure. One line per node.
% =============================================================

\chapter{Background Review}

\section*{Chapter Introduction}
% Three converging literatures (emergence, embodiment, existing systems);
% map to §1.1 dialectics; flag cybernetics thread as background only.

% -------------------------------------------------------------
\section{The Elegance of Emergence}
% Picks up: Bottom-up vs. Top-down.
% Claim: complex global behaviour can arise from simple local rules.

    \subsection{Defining emergence}
    % Emergence as a precise relational property, not a metaphor.

    \subsection{Computational practices that operationalize emergence}
    % CAs, self-organization in biology, engineered swarms / programmable matter.

    \subsection{From rule sets to embodied substrates}
    % Where emergence has stalled in robotics: rigid agents, informational coupling only.
    % Sets up §2.2.

% -------------------------------------------------------------
\section{Matter Does Matter}
% Picks up: Soft vs. Rigid.
% Claim: the body is a computational participant, not a passive platform.

    \subsection{The classical separation of body and controller}
    % Why industrial robotics treats compliance as noise.

    \subsection{Morphological computation: the body as participant}
    % Pfeifer; passive walking; jamming gripper; physical reservoir computing.
    % Müller & Hoffmann's careful boundary.

    \subsection{Morphing matter and programmable self-assembly}
    % Tibbits, programmable matter, smart materials. Bridging vocabulary to §2.3.

    \subsection{Implication: the design question changes}
    % Reframe: "what body shrinks the controller's task?"
    % Caveat: most demos are monolithic; modular is open. Sets up §2.3.

% -------------------------------------------------------------
\section{Robot Unlike Robot}
% Position the thesis against existing robotic systems.
% Title borrowed from Tibbits.

    \subsection{Soft robotics — compliance without comparability}
    % Monolithic; no shared structural primitive across designs.

    \subsection{Modular reconfigurable robotics — comparability without compliance}
    % Comparable units, but rigid; passive material coupling blocked.

    \subsection{At the boundary — soft and modular and reconfigurable}
    % Mochibot, AuxBots, particle robots, Tetraflex, PneuBots.
    % Each pays a price; none asks the thesis question.
    % Anchor: comparison table.

    \subsection{The gap}
    % Soft + modular + topological comparison + locomotion-by-inflation
    % is the thinly populated region. Bridge to Ch.3.
